% English draft based on sections/02_inference_operator_ja.tex
\section{The Inference Operator $F$}

This chapter defines the inference operator used throughout the paper. The operator $F$ is not tied to one implementation such as deduction, induction, Bayesian update, neural generation, or statistical prediction. It denotes the minimal operation by which a state of understanding is advanced into a more justified state.

\subsection{Definition}

Let $\mathcal{X}$ be the state space of inferential states. An inference operator is a map
\[
x_{n+1}=F(x_n),
\]
where $x_n\in\mathcal{X}$ is a provisional state and $F(x_n)$ is the next state generated by inference.

The defining property of $F$ is justification extension. There exists a justification measure
\[
J:\mathcal{X}\to\mathbb{R}
\]
such that, for admissible states,
\[
J(F(x))>J(x).
\]
This measure is not truth. It measures the advance of an inferential chain, not its final correctness.

\subsection{Relation to VED}

In VED terms, inference follows a gradient of closure. It is a local realization of a flow driven by non-closure. A state that appears unresolved generates a direction of further movement. Thus $F$ may be read as the intelligence-layer realization of a closure-gradient flow.

The important consequence is that inference is inherently dynamic. It is an operation for moving, not an operation for terminating movement. If termination is required, it must be supplied by something other than $F$ as such.

\subsection{Branching}

Inference is not necessarily deterministic. In many cases,
\[
F(x_n)=\{x_{n+1}^{(1)},x_{n+1}^{(2)},\ldots\}
\]
is better understood as a branching set of possible continuations. This branching structure connects Part II to the branching explosion described in Part I. The more inference is refined, the more branches can be generated.

For this reason, stronger inference does not automatically mean stronger stability. It may mean that the system has more ways to continue and therefore more ways not to stop.

\subsection{Internal Stopping Predicate}

One might attempt to define a stopping predicate inside inference itself:
\[
S(x)=1 \quad \text{if inference should stop at } x.
\]
But if $S$ is justified by inference, then $S$ requires further inference to justify its own stopping authority. The attempt to define stopping from inside $F$ therefore risks infinite regress. This point is formalized in \S{}3.

\bigskip
